\abstract{РЕФЕРАТ}

Объем работы равен \formbytotal{lastpage}{страниц}{е}{ам}{ам}. Работа содержит \formbytotal{figurecnt}{иллюстраци}{ю}{и}{й} и \arabic{bibcount} библиографических источников. Количество приложений – 1. Фрагменты исходного кода представлены в приложении А.

Перечень ключевых слов: AES-128, симметричное шифрование, Common Lisp, блочное шифрование, PKCS\#7, модуль, функция, библиотека, ключ, тестирование, шифрование, дешифрование, безопасность информации, вектор, блок, алгоритм.

Объектом разработки является библиотека для шифрования данных, реализующая алгоритм AES-128.

Целью курсовой работы является создание программной библиотеки на языке Common Lisp, реализующей алгоритм AES-128 для защиты информации путем шифрования данных с возможностью их последующего дешифрования.

В процессе выполнения работы была спроектирована модульная структура библиотеки, реализованы основные механизмы алгоритма AES. Для обеспечения работы с данными произвольной длины реализована поддержка стандарта дополнения блоков PKCS\#7.

При разработке использовалась среда программирования SBCL (Steel Bank Common Lisp).

Разработанная библиотека успешно прошла тестирование и подтвердила корректность работы.

\selectlanguage{english}
\abstract{ABSTRACT}

The volume of work is \formbytotal{lastpage}{page}{}{s}{s}. The work contains \formbytotal{figurecnt}{illustration}{}{s}{s} and \arabic{bibcount} bibliographic sources. The number of applications is 1. Fragments of the source code are presented in annex A.

List of keywords: AES-128, symmetric encryption, Common Lisp, block cipher, PKCS\#7, module, function, library, key, testing, encryption, decryption, information security, vector, block, algorithm.

The object of development is an encryption library that implements the AES-128 algorithm.

The purpose of the course work is to create a software library in Common Lisp implementing the AES-128 algorithm for information protection via data encryption with the possibility of subsequent decryption.

In the process of work, the modular structure of the library was designed, the main AES mechanisms were implemented. To ensure work with data of arbitrary length, the PKCS\#7 block padding standard was implemented.

The SBCL (Steel Bank Common Lisp) programming environment was used during development.

The developed library was successfully tested and confirmed correct operation.

\selectlanguage{russian}